% Consolidated guide-selftest11 from e02350fe8813d754ca08fbf9437fb39dd83d2210; source fragments preserved.
% 2026-11-24 八批严格二校第8批：历史片段哈希对应修订前版本。

% === BEGIN guide-selftest11-choices.tex blob 080e6872d4036d71cb766929b060411f0980a36f ===
% 《大学物理学学习指导》第5版，自测题11“量子物理”
% 选择11.1.1—11.1.12，印刷230—231，PDF物理238—239页。
\subsection{自测题11：光电、康普顿、物质波与量子态选择题}

\sourcetest{11.1.1}{230}{238}
\problemtext 金属光电逸出需要能量$eU_0$，
单色入射光真空波长$\lambda$必须满足何种阈值条件？
\workings
假定$U_0>0$且为逸出功除以正元电荷的等效电压，在理想单光子模型下发生外光电效应的能量阈值条件是
$hc/\lambda\ge eU_0$，于是
\[
\boxed{\lambda\le\frac{hc}{eU_0}}.
\]
选\textbf{A}。
\crosscheck 波长越短光子能量越大；
不应把阈值不等式方向反写；等号对应最大初动能为零的理想阈值，不保证实际可测光电流非零。

\sourcetest{11.1.2}{230}{238}
\problemtext 同一金属、同一实验中，
照射波长由400nm减至300nm，
问遏止电压如何变化。
\workings
爱因斯坦光电方程$eU_a=hc/\lambda-W$，
假设两种波长都能产生光电子，且遏止电压以非负大小定义，同种金属作差消去逸出功，
\[
\Delta U_a=\frac{hc}{e}
\left(\frac1{300\,\mathrm{nm}}-\frac1{400\,\mathrm{nm}}\right)
\simeq\boxed{+1.033\,\mathrm V}.
\]
故遏止电压\textbf{大小增大约1.03 V}，选\textbf{D}。
\crosscheck 题给选项按约两位小数近似；
不是随波长缩短而降低遏止电压的大小。若记录带符号的负偏压，读数反而减小。

\sourcetest{11.1.3}{230}{238}
\problemtext 保持真空光电管阳极—阴极电势差不变，
仅增大照射同一阴极的同频单色光强，
判断逸出电子的最大初动能及到阳极时最大动能变化。
\workings
光强增大增加单位时间入射光子数与可能光电子数，
但每个光子能量$h\nu$保持不变，
$\boxed{K_{0,\max}=h\nu-W}$不变。
若定义$\Delta V=V_A-V_K$，电子由阴极到阳极时电场功为$e\Delta V$（可正可负），故能到达阳极的最大动能为$K_{A,\max}=h\nu-W+e\Delta V$。若右式为负则没有电子能到达；采用单光子、无空间电荷等影响的理想模型时，改变光强不改变上述能量。
选\textbf{D（均不变）}。
\crosscheck 电流大小可能增大，
不能将光电子数增加误当作单个光电子能量增加。

\sourcetest{11.1.4}{230}{238}
\problemtext X射线散射实验中的散射光可能包含哪些波长成分，
康普顿波长改变是否依赖材料？
\workings
可出现与入射波长相同的弹性（相干）散射分量，
以及波长增长的非弹性康普顿散射分量。
理想自由电子的康普顿位移
\[
\boxed{\Delta\lambda=\frac{h}{m_ec}(1-\cos\theta)}
\]
只依散射角$\theta$和电子静质量，
不因材料的化学名称而改变该基本位移关系。
选\textbf{B}。
\crosscheck 各散射分量的\emph{强度}会明显依赖材料、
束缚结构和入射能量；
不能把“波长位移的理想公式与材料无关”
误写成全部散射现象与材料无关。

\sourcetest{11.1.5}{230}{238}
\problemtext 康普顿散射中散射后光子波长是原来的1.2倍，
求散射后光子能量与反冲电子动能之比。
\workings
光子能量反比波长：
$E'_\gamma=E_\gamma/1.2=5E_\gamma/6$，
自由静止电子反冲动能
$K=E_\gamma-E'_\gamma=E_\gamma/6$，故
\[
\boxed{E'_\gamma/K=5}.
\]
选\textbf{D}。
\crosscheck 题问的是\emph{散射后}光子能量，
不是初始能量与反冲动能之比6。

\sourcetest{11.1.6}{230--231}{238--239}
\problemtext 康普顿散射中原先静止电子被反冲到$0.6c$，
求其动能是静止能量的几倍。
\workings
$\gamma=(1-0.6^2)^{-1/2}=1.25$，
\[
\boxed{\frac{K}{m_ec^2}=\gamma-1=0.25}.
\]
选\textbf{D}。
\crosscheck $0.5m_ev^2=0.18m_ec^2$是经典动能，
在本题速度下误差不可忽略。

\sourcetest{11.1.7}{231}{239}
\problemtext 玻尔氢原子模型中，电子第一与第三轨道速度比$v_1/v_3$。
\workings
玻尔氢原子第$n$轨道角动量量子化
$m_ev_nr_n=n\hbar$，向心库仑力与轨道关系联立
得到$v_n\propto1/n$，故
\[
\boxed{v_1/v_3=3}.
\]
选\textbf{C}。
\crosscheck 此为玻尔圆轨道模型，不能作为现代量子力学
电子真实经典圆轨迹的陈述。

\sourcetest{11.1.8}{231}{239}
\problemtext 氢原子从基态激发到某态需10.19eV；
若又从能量$-0.85$eV的状态跃迁到该激发态，
发射光子能量是多少？
\workings
基态$E_1=-13.6$eV，
加10.19eV得到目标态$E_f=-3.41$eV
（按13.6eV标准近似为$n=2$，$E_2=-3.4$eV）。
从$-0.85$eV（$n=4$）到$n=2$，
采用玻尔能级表发射能量为
\[
\boxed{E_\gamma=(-0.85)-(-3.40)=2.55\,\mathrm{eV}}.
\]
若严格按题给10.19 eV和近似基态-13.6 eV直接代数计算，则$E_f=-3.41$ eV、$E_\gamma=2.56$ eV。两者采用不同的舍入数据约定。
选\textbf{A}。
\crosscheck 原题10.19eV只保留四位数字，
若始终使用它得到2.56eV；
采用精确玻尔13.6eV能级表则是2.55eV，
必须注明数据精度来源。

\sourcetest{11.1.9}{231}{239}
\problemtext 两种不同质量的粒子德布罗意波长相等，
哪一项物理量必然相等？
\workings
德布罗意关系
$\lambda=h/p$给
\[
\boxed{p_1=p_2=h/\lambda},
\]
故是\textbf{动量大小相同}，选\textbf{A}。
\crosscheck 速度、动能和总能量不必相同；
这里无须假定二者具有相同静质量。

\sourcetest{11.1.10}{231}{239}
\problemtext 原图11.1.1显示四幅沿$x$方向的粒子波包(A)—(D)，
问哪种波包有最精确的动量。
\workings
对相同形状、无额外空间相位调制的缩放波包，位置分布越宽，波数分布越窄，故$\Delta p_x=\hbar\,\Delta k$越小；这并非任意量子态间的严格排序定理。
在原图四种相似的局域波包中，
(A)具有最长、最多周期的波列，
波数相对散布最窄，故选\textbf{A}。
\crosscheck \textbf{图形比较为定性：}
原图并未给各波包归一化强度或定量谱宽；
“周期数最多”只是判断较长相干波列动量分布
可能较窄的启发性证据。若需严密排序，
应对四幅波形给出明确解析表达并计算傅里叶谱标准差。

\sourcetest{11.1.11}{231}{239}
\problemtext 波长500nm光子的波长不确定度为$10^{-4}$nm；
按\emph{原题指定的不确定关系}
$\Delta p_x\,\Delta x\ge h$求$x$位置至少多不确定。
\workings
$p=h/\lambda$，仅当波长不确定度为窄带波包内禀谱宽且传播近似沿$x$轴时，小不确定度下$\Delta p_x\simeq h\Delta\lambda/\lambda^2$；单独的仪器测量误差不能确定波包位置宽度。
于是按原题给定的$h$形式，
\[
\boxed{\Delta x\gtrsim\frac h{\Delta p}
=\frac{\lambda^2}{\Delta\lambda}
=\frac{(500\,\mathrm{nm})^2}{10^{-4}\,\mathrm{nm}}
=2.5\,\mathrm m=250\,\mathrm{cm}}.
\]
选\textbf{C}。
\crosscheck \textbf{不等式常数约定：}
在适用的标准差型一维位置—动量模型中，下界为$\hbar/2=h/(4\pi)$，相应估计为$2.5/(4\pi)\simeq0.199$ m。
同时光子的严格位置算符有理论细节；
本条只是在\emph{题目所给波包近似关系}下计算，
不将其当作普适的光子严格定位定理。

\sourcetest{11.1.12}{231}{239}
\problemtext 氢原子L壳层（$n=2$）电子的四量子数
$(n,l,m_l,m_s)$，哪组满足允许值？
\workings
当$n=2$时$l=0,1$；
$l=1$时$m_l=-1,0,1$；
电子$m_s=\pm1/2$。
原选项中满足以上规则的是
\[
\boxed{(n,l,m_l,m_s)=(2,1,-1,1/2)},
\]
选\textbf{B}。
\crosscheck 不能把$n=1$的K壳层或$n=3$的M壳层组合
计入L壳层；$l=0$时$m_l$只能为0。
% === END guide-selftest11-choices.tex ===

% === BEGIN guide-selftest11-blanks.tex blob 73e0d2756638e1dcd468a2ac7dc1908c696ad01c ===
% 《大学物理学学习指导》第5版，自测题11填空11.2.1—11.2.9。
% 源书印刷231—232页，PDF物理239—240页。
\subsection{自测题11：光谱、粒子波长和量子态填空}

\sourcetest{11.2.1}{231}{239}
\problemtext X射线散射后存在两种波长分量，
哪一种称为康普顿散射分量？
\workings
散射光包括与入射波长\textbf{相同}的弹性散射成分，
也包括波长比入射光\textbf{更长}的非弹性散射成分。
后者来自电子反冲的能量转移，称
这就是所填的\textbf{康普顿散射分量}。
\crosscheck 无自由电子反冲时的相干弹性散射不能
因为同在散射光束中就一并称为“康普顿位移”。

\sourcetest{11.2.2}{232}{240}
\problemtext 氢原子巴耳末线系最长波长与最短波长的比值。
\workings
终态$n=2$。
最长线为$3\to2$，
有限束缚态跃迁波长没有取得最小值，其下确界为系列极限$\infty\to2$。
用里德伯关系
\[
\boxed{\frac{\lambda_{\max}}{\lambda_{\min}}
=\frac{1/4}{1/4-1/9}
=\frac95=1.80}.
\]
\crosscheck 最短波长是\emph{巴耳末极限}，
不是从有限$n=4$的次长线误推得到的27/20。

\sourcetest{11.2.3}{232}{240}
\problemtext 氢原子电子位于$n=3$时，
完全电离到自由电子所需能量是多少eV？
\workings
玻尔氢原子能级$E_n=-13.6/n^2$eV，
电离极限$E_\infty=0$，
\[
\boxed{E_{\rm ion}=0-E_3=\frac{13.6}{9}
\simeq1.51\,\mathrm{eV}}.
\]
\crosscheck 不可按基态电离能13.6eV直接填写，
因为原电子已经在$n=3$激发态。

\sourcetest{11.2.4}{232}{240}
\problemtext 中子德布罗意波长$0.2$nm，
取$h=6.63\times10^{-34}$J\,s、
$m=1.67\times10^{-27}$kg，
求其非相对论动能。
\workings
$p=h/\lambda$，$K=p^2/(2m)$，故
\[
\boxed{K=\frac{h^2}{2m\lambda^2}
=\frac{(6.63\times10^{-34})^2}
{2(1.67\times10^{-27})(0.2\times10^{-9})^2}
\simeq3.29\times10^{-21}\,\mathrm J}.
\]
折合约$0.0205$eV。
\crosscheck 中子的静止能量约939MeV，
本题动能极低，非相对论动能公式充分适用。

\sourcetest{11.2.5}{232}{240}
\problemtext 电子局域于长度$\Delta x=0.05$nm范围，
按原书指定$\Delta x\Delta p_x\ge h$估算
$x$方向动量不确定度。
\workings
原书此题明确采用$h$而非$\hbar/2$，
\[
\boxed{\Delta p_x\gtrsim
\frac{6.63\times10^{-34}}
{0.05\times10^{-9}}
\simeq1.33\times10^{-23}\,\mathrm{kg\,m/s}}.
\]
\crosscheck 若用标准差形式
$\Delta x\Delta p_x\ge\hbar/2$，
若$0.05$ nm指位置标准差，下界应再除以$4\pi$，约$1.05\times10^{-24}$ kg\,m/s；若指支持区间全宽$L$，则$\sigma_x\le L/2$，可得$\sigma_p\ge\hbar/L\simeq2.11\times10^{-24}$ kg\,m/s；
两种约定的结果不应混称为同一个精确界限。

\sourcetest{11.2.6}{232}{240}
\problemtext 补全三组$n=2$单电子四量子数中的合法缺项：
(1) $m_l=-1,m_s=-1/2$；
(2) $l=0,m_s=+1/2$；
(3) $l=1,m_l=0$。
\workings
量子数允许范围
$l=0,\ldots,n-1$，
$m_l=-l,\ldots,l$，
$m_s=\pm1/2$。
因此
\[
\boxed{\text{(1)}\ l=1},\qquad
\boxed{\text{(2)}\ m_l=0},\qquad
\boxed{\text{(3)}\ m_s=\pm\frac12}.
\]
\crosscheck $m_l=-1$要求$l\ge1$，
与$n=2$一起唯一确定$l=1$。
原题(3)只给出一个没有其他限制的电子，
允许两种自旋投影而非只能取某一符号。

\sourcetest{11.2.7}{232}{240}
\problemtext 主量子数$n=4$时角量子数$l$
及磁量子数$m_l$可能取值。
\workings
$n=4$时$l=0,1,2,3$。
对每个$l$都有$m_l=-l,\ldots,l$，
综合所有允许$l$，
\[
\boxed{l=0,1,2,3},\qquad
\boxed{m_l=-3,-2,-1,0,1,2,3}.
\]
\crosscheck 虽然总集合含7个$m_l$值，
但并不是每个$l$均可同时取这7个值；
必须满足$|m_l|\le l$。

\sourcetest{11.2.8}{232}{240}
\problemtext 基态氦原子的两名电子用四量子数
$(n,l,m_l,m_s)$分别如何表示？
\workings
在独立电子组态近似下，基态氦的两个电子占据$1s$轨道；
泡利不相容原理要求自旋投影相反，
所以两组为
\[
\boxed{(1,0,0,+1/2)},\qquad
\boxed{(1,0,0,-1/2)}.
\]
\crosscheck 相同的前三量子数是允许的；
不允许两个电子的\emph{四个}量子数同时完全相同。严格两电子基态为空间对称、旋量反对称的单重态，不能将不可分辨电子分别赋予固定身份和固定自旋。

\sourcetest{11.2.9}{232}{240}
\problemtext 原子序数$Z=6$的碳及$Z=14$的硅，
基态电子组态分别是什么？
\workings
在原书采用的独立电子填充近似下，
按照能量顺序和泡利不相容原理，
\[
\boxed{\mathrm C:\ 1s^2\,2s^2\,2p^2},
\]
\[
\boxed{\mathrm{Si}:\ 1s^2\,2s^2\,2p^6\,3s^2\,3p^2}.
\]
\crosscheck 电子占据数分别为$2+2+2=6$，
及$2+2+6+2+2=14$；
开放$p$亚层的具体自旋排列还受洪特规则影响。
% === END guide-selftest11-blanks.tex ===

% === BEGIN guide-selftest11-solutions.tex blob 111621b735b97c8ebd2e2e9faf8e26aef5b80260 ===
% 《大学物理学学习指导》自测题11计算11.3.1—11.3.6
% 原书印刷232，PDF物理240页。源题数据、单位与边界逐项核对。
\subsection{自测题11：光电、康普顿、氢谱和势阱计算题}

\sourcetest{11.3.1}{232}{240}
\problemtext 410nm单色光照射金属，最大光电子动能为1.0eV，
求能使该金属发生光电效应的单色光最大波长。
\workings
由爱因斯坦光电方程
$W=hc/\lambda-K_{\max}$，采用
$hc\simeq1239.842\,\mathrm{eV\,nm}$，
\[
 W=\frac{1239.842}{410}-1.0
 \simeq2.024\,\mathrm{eV}.
\]
红限波长由$hc/\lambda_0=W$给出，
\[
\boxed{\lambda_{\max}=\lambda_0
=\frac{1239.842}{1239.842/410-1}\,\mathrm{nm}
\simeq612.6\,\mathrm{nm}}.
\]
\crosscheck 红限波长应\emph{长于}410nm，
因为原入射光子有1eV的多余动能可供电子带走。

\sourcetest{11.3.2}{232}{240}
\problemtext 0.0500nm的X射线在自由静止电子上发生康普顿散射；
散射后波长0.0522nm。
求光子的散射角及电子得到的动能。
\workings
康普顿波长
$\lambda_C=h/(m_ec)=0.00242631$nm。
波长增加
$\Delta\lambda=0.0522-0.0500=0.0022$nm，
\[
\cos\theta=1-\frac{\Delta\lambda}{\lambda_C}
=1-\frac{0.0022}{0.00242631}
\simeq0.09328,
\]
因此
\[
\boxed{\theta\simeq84.65^\circ}.
\]
电子初始静止，反冲动能为光子损失的能量：
\[
\begin{aligned}
K_e&=hc\left(\frac1{\lambda_0}-\frac1{\lambda'}\right)\\
&\simeq(1239.842\,\mathrm{eV\,nm})
\left(\frac1{0.0500}-\frac1{0.0522}\right).
\end{aligned}
\]
从而
\[
\boxed{K_e\simeq1.045\times10^3\,\mathrm{eV}
=1.045\,\mathrm{keV}}.
\]
\crosscheck $\lambda'>\lambda_0$，故电子得到正动能；
将$\lambda,\lambda'$倒置会产生负能量，违反题设。

\sourcetest{11.3.3}{232}{240}
\problemtext 氢原子电子从基态跃迁至主量子数$n$的激发态，
这时玻尔轨道半径变为原来9倍，求原子从该态
直接返回基态时发射光子的频率。
\workings
玻尔轨道半径$r_n=n^2a_0$，
$r_n/r_1=9$给$n=3$。
氢能级$E_n=-13.6/n^2$eV，
$n=3\to1$光子能量
\[
\Delta E=13.6\left(1-\frac1{9}\right)
\simeq12.0889\,\mathrm{eV}.
\]
因此
\[
\boxed{\nu=\frac{\Delta E}{h}
=\frac{12.0889\,e}{h}
\simeq2.923\times10^{15}\,\mathrm{Hz}}.
\]
\crosscheck 上述直接单光子跃迁还要求满足实际角动量选择定则（例如电偶极$3p\to1s$允许，非所有$n=3$子态均可直接跃迁）；题中“基态跃迁到激发态”本身是\emph{吸收}能量；
要谈\emph{发射}光子需明确从$n=3$返回较低能级，
此处按源题的直接返回基态模型给出频率。

\sourcetest{11.3.4}{232}{240}
\problemtext 氢原子从某个初始态跃迁到
\emph{从基态激发需10.19eV的能级}时
发射波长486nm的光子，
求初始态的能量与主量子数。
\workings
从基态$E_1=-13.6$eV激发10.19eV，
终态能量约$E_f=-3.41$eV，
即按玻尔能级取$n_f=2$、$E_f\simeq-3.4$eV。
光子能量
$E_\gamma=hc/(486\,\mathrm{nm})\simeq2.551$eV。
于是初始电子能量
\[
E_i=E_f+E_\gamma\simeq-3.4+2.55
=-0.85\,\mathrm{eV}.
\]
解$E_i=-13.6/n_i^2$，
\[
\boxed{n_i=4},\qquad
\boxed{E_i=-0.85\,\mathrm{eV}}.
\]
\crosscheck 巴耳末线486nm对应$4\to2$；
不能因为题干出现“从基态激发能10.19eV”
就错把486nm理解成$1\to2$发射。

\sourcetest{11.3.5}{232}{240}
\problemtext 若采用能量等效质量$m_{\rm eq}=E/c^2$的旧教材表述，
求波长550nm光子的等效质量（kg）。
\workings
光子能量$E=hc/\lambda$，所以
\[
\boxed{m_{\rm eq}=\frac{E}{c^2}
=\frac{h}{\lambda c}
=\frac{6.62607015\times10^{-34}}
{(550\times10^{-9})(299792458)}
\simeq4.02\times10^{-36}\,\mathrm{kg}}.
\]
\crosscheck 这不是光子的\emph{静质量}；
现代相对论表述中光子静质量为零，
这里的$m_{\rm eq}$仅是光子能量除以$c^2$。

\sourcetest{11.3.6}{232}{240}
\problemtext 一维无限深势阱$0\le x\le a$内，
已归一化波函数为
$\psi(x)=\sqrt{2/a}\sin(\pi x/a)$，
求找到粒子概率密度最大的位置。
\workings
概率密度
\[
\rho(x)=|\psi(x)|^2
=\frac2a\sin^2\left(\frac{\pi x}{a}\right).
\]
在$[0,a]$范围内，
$\sin^2(\pi x/a)$于$\pi x/a=\pi/2$处达到最大1，
故
\[
\boxed{x=\frac a2},\qquad
\boxed{\rho_{\max}=\frac2a}.
\]
\crosscheck 势阱两端$x=0,a$的波函数为零，
故概率密度最大的位置在阱的正中心，而不是两端。
% === END guide-selftest11-solutions.tex ===
